\documentclass[12pt]{article}
\usepackage[margin=1in]{geometry}
\usepackage{amsmath,amssymb,booktabs,array,threeparttable}

\begin{document}

\providecommand{\StructuralRobustnessPath}{appendix/structural-robustness}
\providecommand{\E}{\mathbb{E}}
\providecommand{\Cov}{\operatorname{Cov}}

\subsection{Which Data Discipline the Structural Objects?}
\label{app:structural-object-sensitivity}

Our structural conclusions concern equilibrium objects rather than any one
primitive parameter in isolation. In particular, the model maps the joint
posterior into four economically meaningful objects: the social multiplier,
the social-image return, the private value of the incentives, and the
visibility of deworming. Table~\ref{tab:object-sensitivity} reports a Bayesian
likelihood-sensitivity diagnostic showing which parts of the data discipline
each object.

Let $Q(\theta)$ denote a structural object and let $\ell_b(\theta)$ be the log
likelihood for data block $b$. If that block is assigned weight $\omega_b$,
the local response of its posterior mean is
\[
  \left.
  \frac{\partial \E_{\omega}[Q(\theta)\mid Y]}
       {\partial \omega_b}
  \right|_{\omega_b=1}
  =
  \Cov\!\left(Q(\theta),\ell_b(\theta)\mid Y\right).
\]
We divide this covariance by the posterior standard deviation of $Q$, so an
entry measures the movement in the posterior mean, in posterior-standard-
deviation units, induced by a local increase in the weight placed on a data
block. Table~\ref{tab:object-sensitivity} evaluates the structural schedules
at 1.5 km and reports each treatment arm separately. The first row within each
group gives the root mean square of the four arm-specific sensitivities.

\begin{table}[htbp]
  \centering
  \small
  \caption{Absolute posterior-mean sensitivity of structural objects to the data}
  \label{tab:object-sensitivity}
  
\begin{tabular}[t]{>{\raggedright\arraybackslash}p{5.4cm}rrr}
\toprule
Structural object & Deworming take-up & Observability & WTP\\
\midrule
\em{Social multiplier, $\tilde{SM}(\textrm{treat}_z,d)$} & \textbf{0.32} & 0.23 & 0.01\\
\hspace{1em}Control & \textbf{0.50} & 0.09 & 0.00\\
\hspace{1em}Ink & \textbf{0.32} & 0.11 & 0.01\\
\hspace{1em}Calendar & \textbf{0.21} & 0.03 & 0.00\\
\hspace{1em}Bracelet & 0.13 & \textbf{0.43} & 0.02\\
\midrule
\em{Social-image return, $\mu(\textrm{treat}_z,d)\Delta(w^\ast(\textrm{treat}_z,d))$} & \textbf{0.50} & 0.11 & 0.00\\
\hspace{1em}Control & \textbf{0.49} & 0.11 & 0.00\\
\hspace{1em}Ink & \textbf{0.51} & 0.11 & 0.00\\
\hspace{1em}Calendar & \textbf{0.49} & 0.11 & 0.00\\
\hspace{1em}Bracelet & \textbf{0.51} & 0.12 & 0.00\\
\midrule
\em{Private valuation (Calendar relative to Bracelet), $\psi$} & \textbf{0.12} & 0.07 & 0.09\\
\midrule
\em{Observability, $p_{\textrm{Observed}}(\textrm{treat}_z,d)$} & 0.06 & \textbf{0.11} & 0.03\\
\hspace{1em}Control & 0.07 & \textbf{0.08} & 0.05\\
\hspace{1em}Ink & 0.04 & \textbf{0.08} & 0.03\\
\hspace{1em}Calendar & 0.06 & \textbf{0.12} & 0.01\\
\hspace{1em}Bracelet & 0.07 & \textbf{0.17} & 0.01\\
\bottomrule
\end{tabular}

  \begin{minipage}{0.90\linewidth}
    \vspace{0.4em}\footnotesize
    \textit{Notes:} Entries are absolute standardized posterior-mean
    sensitivities. Larger entries indicate that the posterior mean is more
    responsive to the likelihood block. For arm-specific objects, the
    italicized group row reports the RMS across arms. Bold denotes the largest
    sensitivity in each row. All schedules are evaluated at 1.5 km.
  \end{minipage}
\end{table}

The pattern mirrors the study design. Experimental take-up across arms and
distances anchors the magnitude of the social-image return. The multiplier is
jointly disciplined by take-up and observability beliefs: take-up is
especially important for Control, whereas observability beliefs are especially
important for Bracelet. Belief data directly discipline predicted visibility,
and the randomized WTP experiment isolates private incentive values from
signaling value. Thus distinct experimental and survey modules discipline the
combinations of primitives that enter the decomposition and policy
counterfactuals.

\end{document}
